\documentclass[11pt,a4paper]{article}

% ---------- arXiv-compatible packages ----------
\usepackage[utf8]{inputenc}
\usepackage[T1]{fontenc}
\usepackage{hyperref}
\usepackage{amsmath,amssymb,amsthm}
\usepackage{mathtools}
\usepackage{booktabs}
\usepackage{graphicx}
\usepackage{xcolor}
\usepackage{algorithm}
\usepackage{algpseudocode}
\usepackage{geometry}
\usepackage{microtype}
% enumitem not available — using basic LaTeX lists instead
% caption/subcaption not available — using basic figure handling

% Fix hyperref destination conflict with algorithm float
\makeatletter
\providecommand*{\theHalgorithm}{\thealgorithm}
\makeatother

\geometry{margin=1in}

% ---------- metadata for arXiv ----------
\hypersetup{
  colorlinks=true,
  linkcolor=blue,
  citecolor=blue,
  urlcolor=blue,
  pdftitle={UraionSpec: A Faithful Implementation of Confidence-Scheduled Speculative Decoding},
  pdfauthor={Uraion Labs},
  pdfkeywords={speculative decoding, DSpark, semi-autoregressive, confidence scheduling, LLM inference}
}

% ---------- convenience macros ----------
\newcommand{\UraionSpec}{\textsc{UraionSpec}}
\newcommand{\DSpark}{\textsc{DSpark}}
\newcommand{\DFlash}{\textsc{DFlash}}
\newcommand{\DeepSpec}{\textsc{DeepSpec}}
\newcommand{\Eagle}{\textsc{Eagle}}
\newcommand{\MTP}{\textsc{MTP}}

\DeclareMathOperator{\softmax}{softmax}
\DeclareMathOperator{\sigmoid}{sigmoid}
\DeclareMathOperator{\CumProd}{CumProd}
\DeclareMathOperator{\SPS}{SPS}

\begin{document}

% ======================================================================
\title{\textbf{\UraionSpec}: A Faithful Implementation of \\ Confidence-Scheduled Speculative Decoding}

\author{%
  \normalsize Uraion Labs \\[4pt]
  \texttt{https://uraionlabs.com} \\[2pt]
  \texttt{https://huggingface.co/UraionLabs/UraionSpec}
}

\date{\normalsize Technical Report -- ICML 2026 Systems Track}

\maketitle

% ======================================================================
\begin{abstract}
Speculative decoding accelerates autoregressive language model inference by
generating multiple candidate tokens from a lightweight draft model, then
verifying them in parallel with the target model. \DSpark (Cheng et al.,~2026)
introduces two key innovations: semi-autoregressive generation via a parallel
backbone with a lightweight sequential head, and confidence-scheduled
verification that dynamically adjusts verification length based on predicted
acceptance probabilities and engine throughput.

We present \UraionSpec, a clean, modular, verified implementation of
\DSpark. It faithfully reproduces the core
components---the Markov and RNN sequential heads, confidence head, Sequential
Temperature Scaling~(STS) calibration, hardware-aware prefix scheduler
(Algorithm~1), and the three-term training objective---and includes a
\DFlash-style parallel backbone with target model KV injection as a
standalone module for integration into the draft model pipeline.
With 80 passing unit tests, a fully
configurable training pipeline, comprehensive evaluation tools, and
end-to-end smoke tests, \UraionSpec is designed for reproducible research
and practical experimentation. The codebase is publicly available on
GitHub and HuggingFace under the MIT license.
\end{abstract}

% ======================================================================
\section{Introduction}
\label{sec:introduction}

The rapid growth of large language models (LLMs) has made inference latency a
critical bottleneck in production deployments. Autoregressive decoding ---
generating one token at a time with sequential dependence on all previous
tokens --- is inherently bandwidth-bound on modern hardware because the
attention mechanism requires loading the full model parameters from memory
for each generated token~\cite{pope2023efficiently, leviathan2023fast}.

Speculative decoding~\cite{leviathan2023fast, chen2023accelerating} breaks
this bottleneck by introducing a two-phase process: a lightweight draft model
proposes multiple candidate tokens in one forward pass, and the target model
verifies them in parallel, yielding a theoretical speedup proportional to the
acceptance rate. This approach preserves the exact output distribution of the
target model, making it a lossless acceleration technique.

The draft model design determines the fundamental trade-off: more accurate
draft models produce longer accepted prefixes but incur higher drafting cost,
while simpler draft models are cheaper but suffer from low acceptance rates.
\DSpark~\cite{cheng2026dspark} addresses this trade-off through two key
contributions:

\begin{enumerate}
  \item \textbf{Semi-autoregressive generation}: A parallel backbone
    (the \DFlash architecture) processes all proposal positions in a
    single forward pass, while a lightweight sequential head (Markov or RNN)
    injects inter-token dependency. This combines the speed of parallel
    drafters~\cite{stern2018blockwise, xia2024medusa} with the quality of
    autoregressive ones~\cite{miao2025eagle3}.
  \item \textbf{Confidence-scheduled verification}: A trained confidence head
    predicts per-position acceptance probabilities, and a hardware-aware
    scheduler dynamically selects verification lengths to maximize expected
    throughput. This prevents wasting target-model compute on low-confidence
    suffix tokens under heavy load.
\end{enumerate}

\paragraph{Our contribution.}
We present \UraionSpec, a faithful, modular, and verified implementation of
\DSpark. Our implementation includes all core components listed in the
original paper, and a standalone \DFlash-style parallel backbone with target model KV
injection module, available for future integration into the draft model pipeline.
The codebase is designed for clarity and reproducibility rather than
production scale, making it suitable for academic research, ablation studies,
and educational purposes. Key features include:

\begin{itemize}
  \item Complete implementation of all \DSpark components (80 unit tests)
  \item \DFlash-style backbone module with target model KV injection
  \item Markov, gated Markov, and RNN sequential heads
  \item Confidence head with analytical acceptance-rate supervision
  \item Sequential Temperature Scaling (STS) calibration
  \item Hardware-aware prefix scheduler (Algorithm~1)
  \item Training pipeline with three-term objective
  \item Comprehensive evaluation and benchmarking tools
  \item All code under MIT license on GitHub and HuggingFace
\end{itemize}

% ======================================================================
\section{Background and Related Work}
\label{sec:background}

\subsection{Speculative Decoding}
Standard speculative decoding~\cite{leviathan2023fast,
chen2023accelerating} works as follows: given a target model $p_t$ and a
draft model $p_d$, at each decoding step the draft model proposes $\gamma$
candidate tokens $x_1, \ldots, x_\gamma$ left-to-right. The target model
then verifies all $\gamma$ tokens in a single forward pass. Each token $x_k$
is accepted with probability
\begin{equation}
  \textstyle
  P(\text{accept } x_k) = \min\left(1,\; \frac{p_t(x_k \mid x_{<k})}{p_d(x_k \mid x_{<k})}\right),
  \label{eq:acceptance}
\end{equation}
and the first rejection marks the end of the accepted prefix. A bonus token
is then sampled from the target model's residual distribution at the rejection
position, ensuring the final distribution matches the target model exactly.

The expected number of accepted tokens $\tau$ is
\begin{equation}
  \tau = 1 + \sum_{j=1}^\gamma \prod_{i=1}^j \alpha_i,
  \qquad
  \alpha_i = 1 - \frac12 \|p_t(\cdot) - p_d(\cdot)\|_1,
  \label{eq:expected_tau}
\end{equation}
where $\alpha_i$ is the per-position acceptance probability. The throughput
benefit comes from the fact that verifying $\gamma$ tokens costs roughly the
same as generating one token autoregressively, so effective speedup is
$\tau / 1$ (plus overhead).

\subsection{Draft Model Architectures}
Draft model architectures span a spectrum from simple heuristics to complex
autoregressive networks:

\paragraph{Independent drafters.}
Early work~\cite{stern2018blockwise} proposed independent drafters that
predict each position without conditioning on other draft tokens. While
fast (a single forward pass suffices), these suffer from ``multi-modal
collision'' where later tokens become incoherent with earlier ones.

\paragraph{Autoregressive drafters.}
\Eagle~\cite{miao2025eagle3} and similar methods use a full autoregressive
draft model, achieving high acceptance rates but requiring $\gamma$
sequential forward passes, offsetting the parallelism benefit.

\paragraph{Hybrid approaches.}
\DSpark's semi-autoregressive design splits the difference: a parallel
backbone handles bulk computation, while a tiny sequential head (Markov or
RNN) adds token-to-token conditioning at negligible cost.

\subsection{Confidence-Based Verification}
Standard speculative decoding verifies all $\gamma$ draft tokens regardless
of their quality. Under load, this wastes target-model batch capacity on
tokens that are likely to be rejected. \DSpark introduces a learned
confidence head $c_k = \sigma(w^T [h_k; W_1 x_{k-1}])$ that predicts the
conditional acceptance probability at each position, and a hardware-aware
scheduler that selects per-request verification lengths to maximize
throughput $\Theta = \tau \cdot \SPS(B)$, where $\SPS(B)$ is the
pre-profiled steps-per-second curve of the inference engine at batch size
$B$.

\subsection{Sequential Temperature Scaling}
Calibration of confidence scores is important for the scheduler's
performance. \DSpark proposes Sequential Temperature Scaling (STS), which
calibrates confidence scores left-to-right. For each position $k$, a
temperature $T_k$ is found that minimizes the Expected Calibration Error
(ECE) of the cumulative survival probability $\prod_{i=1}^k c_i$, while
keeping already-calibrated positions $<k$ fixed. The calibrated score at
position $k$ is
\begin{equation}
  \tilde{c}_k = c_k^{1/T_k}.
\end{equation}
Because temperature scaling is order-preserving, it corrects probability
magnitudes without disrupting relative rankings.

% ======================================================================
\section{\UraionSpec: Architecture and Implementation}
\label{sec:method}

\UraionSpec is organized as a modular Python package with clear separation
between model components, decoding logic, training infrastructure,
calibration, and evaluation. Figure~\ref{fig:architecture} shows the
overall structure.

\begin{figure}[t]
  \centering
  \begin{minipage}{0.85\textwidth}
  \ttfamily\footnotesize
  \begin{verbatim}
UraionSpec/
  src/uraionspec/
    models/          -- Draft model components
      markov_head.py       -- Low-rank transition bias
      rnn_head.py          -- GRU (Gated Recurrent Unit)-like recurrent head
      confidence_head.py   -- Per-position acceptance predictor
      dflash_backbone.py   -- DFlash backbone w/ KV injection
      draft_model.py       -- Combined DSpark draft model
    decoding/
      acceptance.py        -- Lossless rejection sampling
      scheduler.py         -- Algorithm 1: Hardware-aware scheduler
      speculative.py       -- Draft -> verify -> accept loop
    training/
      dataset.py           -- Anchor-block dataset preparation
      losses.py            -- CE + TV + confidence (Eq. 12)
      train_drafter.py     -- Training loop (frozen target)
      cache_targets.py     -- Target logit cache generation
    calibration/
      sts.py               -- Sequential Temperature Scaling
    evaluation/
      eval_acceptance.py   -- Acceptance rate/length metrics
      benchmark_latency.py -- Vanilla vs speculative latency
  \end{verbatim}
  \end{minipage}
  \caption{\UraionSpec package structure. Each module is independently
  testable and documented.}
  \label{fig:architecture}
\end{figure}

\subsection{Sequential Heads}
\label{subsec:sequential-heads}

\UraionSpec implements three sequential head variants, each with
different trade-offs between expressiveness and computational cost.

\paragraph{Vanilla Markov head.}
The simplest sequential head implements a low-rank transition bias:
\begin{equation}
  B(x_{k-1}, x_k) = W_1[x_{k-1}] \cdot W_2,
  \qquad
  W_1 \in \mathbb{R}^{V \times r},\; W_2 \in \mathbb{R}^{r \times V},
\end{equation}
where $r = 256$ is the rank and $V$ is the vocabulary size. This head is
memoryless: the bias depends only on the immediately preceding token, not
on full prefix history.

\paragraph{Gated Markov head.}
The gated variant modulates the Markov bias using the backbone hidden state:
\begin{equation}
  g_k = \sigma\bigl(W_g [h_k ;\; W_1[x_{k-1}]]\bigr), \qquad
  B_k = W_2\bigl(g_k \odot W_1[x_{k-1}]\bigr).
\end{equation}
This allows the backbone representation to influence how much Markov bias to
apply, providing context-dependent transition modeling.

\paragraph{RNN head.}
For full prefix history, we implement a GRU (Gated Recurrent Unit)-like recurrent head:
\begin{align}
  z_k &= [s_{k-1} ;\; W_1[x_{k-1}] ;\; h_k], \\
  s_k &= \sigma(W_g z_k) \odot s_{k-1} + (1 - \sigma(W_g z_k)) \odot \tanh(W_c z_k), \\
  B_k &= W_2^{\mathsf{T}} \tanh(W_o z_k).
\end{align}
The recurrent state $s_k$ accumulates information across the entire prefix
$x_{<k}$, enabling the draft model to capture longer-range dependencies
than the Markov variants.

\subsection{Confidence Head}

The confidence head predicts the per-position conditional acceptance
probability:
\begin{equation}
  c_k = \sigmoid\bigl(w^{\mathsf{T}} [h_k ;\; W_1[x_{k-1}]]\bigr),
\end{equation}
where $h_k$ is the backbone hidden state at position $k$ and $W_1[x_{k-1}]$
is the Markov embedding of the previous token. The head is supervised by the
analytical acceptance rate:
\begin{equation}
  c_k^* = 1 - \tfrac12 \|p_d(\cdot \mid x_{<k}) - p_t(\cdot \mid x_{<k})\|_1,
\end{equation}
which is the per-position conditional probability that $x_k$ is accepted
under the distribution-matching formulation of speculative decoding.

\subsection{DFlash-Style Backbone with KV Injection}
\label{subsec:backbone}

The parallel backbone processes all $\gamma$ proposal positions in a single
forward pass. A key design choice in \DSpark is that each backbone layer
has access to the target model's hidden representations at the corresponding
layer, which we implement through KV injection.

\paragraph{DFlashAttention.}
Each attention layer constructs its key-value pairs by concatenating target
context features with draft features:
\begin{align}
  K &= [W_K \cdot H^\text{ctx} ;\; W_K \cdot H^\text{draft}], \\
  V &= [W_V \cdot H^\text{ctx} ;\; W_V \cdot H^\text{draft}],
\end{align}
where $H^\text{ctx}$ is the target model's hidden states at the same layer
and $H^\text{draft}$ is the draft model's hidden states. The query is
computed from draft hidden states only:
\begin{equation}
  Q = W_Q \cdot H^\text{draft},
\end{equation}
and attention proceeds as standard scaled dot-product attention over the
concatenated KV sequence:
\begin{equation}
  \text{Attention}(Q, K, V) = \softmax\left(\frac{Q K^{\mathsf{T}}}{\sqrt{d_k}}\right) V.
\end{equation}
This allows draft representations to condition on the rich contextual
representations from the (frozen) target model.

\paragraph{Block-diagonal attention mask.}
During multi-block training, each block of $\gamma$ draft tokens attends
bidirectionally to all context tokens and to other draft tokens within the
same block, but is masked from attending to draft tokens in other blocks:
\begin{equation}
  M(i,j) = \begin{cases}
    0 &\text{if token $i$ is in block $b_i$, token $j$ is in context or block $b_i$},\\
    -\infty &\text{otherwise}.
  \end{cases}
\end{equation}

\subsection{Training Objective}
\label{subsec:training}

The training objective follows the \DSpark paper (Equation~12):
\begin{equation}
  \mathcal{L} = \alpha_\text{CE} \mathcal{L}_\text{CE} + \alpha_\text{TV} \mathcal{L}_\text{TV} + \alpha_\text{conf} \mathcal{L}_\text{conf},
\end{equation}
with default coefficients $\alpha_\text{CE}=0.1$, $\alpha_\text{TV}=0.9$,
$\alpha_\text{conf}=1.0$. All terms are position-weighted by
$w_k = e^{-(k-1)/\gamma}$ to emphasize earlier positions where acceptance
has a larger impact on overall throughput.

\paragraph{Cross-entropy loss.}
$\mathcal{L}_\text{CE}$ is the standard next-token prediction loss on the
draft model's output, using teacher-forced ground-truth tokens.

\paragraph{Total variation loss.}
$\mathcal{L}_\text{TV}$ minimizes the L1 distance between draft and target
distributions, which is a proxy for maximizing the acceptance rate:
\begin{equation}
  \mathcal{L}_\text{TV} = \|p_d(\cdot \mid x_{<k}) - p_t(\cdot \mid x_{<k})\|_1.
\end{equation}

\paragraph{Confidence loss.}
$\mathcal{L}_\text{conf}$ is binary cross-entropy between the predicted
confidence $c_k$ and the analytical acceptance rate $c_k^*$:
\begin{equation}
  \mathcal{L}_\text{conf} = \text{BCE}\bigl(c_k,\; c_k^*\bigr).
\end{equation}

\subsection{Hardware-Aware Prefix Scheduler}
\label{subsec:scheduler}

We implement Algorithm~1 from the \DSpark paper for dynamic verification
length selection. Given a set of $R$ concurrent requests with confidence
scores $\{c_{r,1}, \ldots, c_{r,\gamma}\}_{r=1}^R$ and a pre-profiled
throughput curve $\SPS(B)$, the scheduler computes:

\begin{algorithm}[t]
\caption{Hardware-Aware Prefix Scheduler (Algorithm~1)}
\label{alg:scheduler}
\begin{algorithmic}[1]
\State \textbf{Input:} confidence scores $\{c_{r,1}, \ldots, c_{r,\gamma}\}_{r=1}^R$, throughput curve $\SPS(B)$
\State \textbf{Output:} verification lengths $\{\ell_r\}_{r=1}^R$
\For{each request $r$}
  \State $a_{r,j} \gets \prod_{i=1}^j c_{r,i}$ \Comment{Prefix survival probs}
\EndFor
\State $C \gets \{(r,j) \mid a_{r,j} > \varepsilon\}$ \Comment{Valid candidates}
\State Sort $C$ descending by $a_{r,j}$
\State $\ell_r \gets 0,\; B \gets R,\; \tau \gets R$ \Comment{Init: 1 anchor per request}
\State $\Theta^* \gets \tau \cdot \SPS(B)$
\For{$(r,j) \in C$ in sorted order}
  \If{$\ell_r \neq j$} \Comment{Non-contiguous prefix: skip}
    \State \textbf{continue}
  \EndIf
  \State $\ell_r \gets j+1,\; B \gets B+1,\; \tau \gets \tau + a_{r,j}$
  \State $\Theta \gets \tau \cdot \SPS(B)$
  \If{$\Theta > \Theta^*$}
    \State $\Theta^* \gets \Theta$
  \Else
    \State $\ell_r \gets j$ \Comment{Early stop: throughput saturated}
    \State \textbf{break}
  \EndIf
\EndFor
\State \Return $\{\ell_r\}_{r=1}^R$
\end{algorithmic}
\end{algorithm}

The scheduler's key property is monotonicity: it processes candidates in
order of decreasing survival probability, which ensures that extending a
request's verification length only happens when it adds the globally most
valuable token. The early stopping condition guarantees the
non-anticipating property --- the scheduler never extends a request's
length based on survival probabilities of future tokens, which cannot be
known at the current step.

A fallback \texttt{StaticScheduler} with fixed-length and static-threshold
modes is also provided for comparison.

\subsection{Sequential Temperature Scaling}
\label{subsec:sts}

We implement STS calibration as described in Section~3.2.1 of the \DSpark
paper. The algorithm processes positions $k = 1, \ldots, \gamma$ left to
right. At each position $k$:

\begin{enumerate}
  \item Compute the cumulative survival probability using already-calibrated
    scores for positions $<k$ and the raw score at position $k$:
    \[
    \hat{a}_k = \prod_{j=1}^{k-1} c_j^{1/T_j} \cdot c_k.
    \]
  \item Find $T_k$ via grid search over $(0.1, 10.0)$ minimizing ECE
    between $\hat{a}_k$ and the empirical cumulative acceptance labels.
  \item Record $T_k$ for inference-time application.
\end{enumerate}

Our implementation uses 100-point grid search to ensure reliable temperature
selection and 15-bin ECE computation for calibration quality assessment.

% ======================================================================
\section{Experimental Verification}
\label{sec:experiments}

We validate the \UraionSpec implementation through unit tests, gradient
checks, smoke training, and component-level verification. All experiments
were conducted on a MacBookPro12,1 (Intel i5-5257U, 7.8~GB RAM, CPU only)
running Ubuntu~24~LTS, targeting the \text{Qwen/Qwen3-0.6B} model.

\subsection{Unit Tests}
\label{subsec:unit-tests}

The codebase includes 80 unit tests spanning all major components. Table
\ref{tab:tests} summarizes test coverage by module.

\begin{table}[t]
\centering
\small
\caption{Unit test coverage by module. All 80 tests pass.}
\label{tab:tests}
\begin{tabular}{lcc}
\toprule
Test Suite & Tests & Coverage \\
\midrule
Acceptance rule ($\min(1,\; p_t/p_d)$) & 8 & Full: all, rejected, partial,\\
 & & bonus token, batch independence,\\
 & & expected length \\
Markov \& RNN heads & 15 & Vanilla, gated, RNN forward shapes,\\
 & & gradients, sampling \\
Scheduler & 14 & Throughput profile, Algorithm~1,\\
 & & static fallback, zero conf. \\
Shape/gradient checks & 10 & Draft model, losses, conf. head,\\
 & & end-to-end cycle \\
STS calibration & 8 & ECE, temperature fitting,\\
 & & fit/transform, calibrator API \\
DFlash backbone & 17 & Attention, decoder layer, full backbone,\\
 & & GQA, masks, gradient flow, edge cases \\
Sampling utilities & 8 & Residual, greedy, temperature, gather \\
\bottomrule
\end{tabular}
\end{table}

\subsection{Component Verification}
\label{subsec:component-verification}

\paragraph{Acceptance rule.}
We verify the lossless speculative decoding acceptance rule
(Eq.~\ref{eq:acceptance}) across multiple scenarios:
\begin{itemize}
  \item \textbf{All accepted}: when $p_t = p_d$, all $\gamma$ tokens are
    accepted with probability~1.
  \item \textbf{First rejection}: the first token with $p_t < p_d$ triggers
    rejection, and all subsequent tokens are discarded.
  \item \textbf{Bonus token}: the bonus token is correctly sampled from the
    residual distribution $\max(0, p_t - p_d) / \| \cdot \|_1$.
  \item \textbf{Batch independence}: acceptance decisions across batch
    elements are independent.
\end{itemize}

\paragraph{DFlash backbone.}
The \DFlash-style backbone passes all shape and gradient checks:
\begin{itemize}
  \item Forward pass produces correct output shapes for single/GQA (Grouped Query Attention)/masked
    attention configurations.
  \item Gradient flow verified through attention, decoder layer, and full
    backbone stack.
  \item Empty draft ($L_d=0$) and single-token draft edge cases handled correctly.
  \item All three activation functions (GELU, ReLU, SiLU) work.
  \item Block-diagonal attention mask correctly enforces intra-block
    attention and cross-block isolation.
\end{itemize}

\paragraph{Training objective.}
The three-term loss (Equation~12 of the \DSpark{} paper) is verified through a smoke training run:
\begin{itemize}
  \item 32 samples from the Capybara dataset, block size $\gamma=4$.
  \item All three loss terms (CE, TV, confidence) are non-negative and
    decrease over training steps.
  \item Position weighting correctly assigns higher weight to earlier
    positions.
  \item Analytical acceptance rate supervision provides valid training
    signal for the confidence head.
\end{itemize}

\subsection{Limitations}
\label{subsec:limitations}

While \UraionSpec faithfully implements the \DSpark algorithm, several
limitations should be noted:

\begin{enumerate}
  \item \textbf{No multi-GPU training}: The training pipeline is
    single-device only. Full-scale reproduction requires multi-GPU
    infrastructure beyond our current scope.
  \item \textbf{Synthetic throughput profile}: The hardware-aware scheduler
    uses a default SPS curve representative of an A100 GPU. For
    production deployment, users should profile their own inference engine
    and pass the real curve.
  \item \textbf{Small-scale validation}: Our experiments are limited to
    CPU-based smoke testing of \text{Qwen3-0.6B}. Full training and
    benchmarking on GPU hardware (e.g., Colab A100) is the natural next
    step.
  \item \textbf{No vLLM integration}: Unlike the official \DeepSpec{}
    repository, \UraionSpec does not integrate with production serving
    frameworks. We prioritize algorithm clarity over production readiness.
\end{enumerate}

% ======================================================================
\section{Reproducibility}
\label{sec:reproducibility}

All code is publicly available under the MIT license:

\begin{itemize}
  \item \textbf{GitHub}: \url{https://github.com/arnavprabhu/UraionSpec}
  \item \textbf{HuggingFace}: \url{https://huggingface.co/UraionLabs/UraionSpec}
\end{itemize}

The repository includes complete documentation, implementation notes, a
reproduction report, and model card templates. To replicate our experiments:
\begin{verbatim}
pip install -e .
pytest tests/ -v

# Smoke training (CPU, 32 samples, 5 steps)
python scripts/smoke_train.py \
    --target Qwen/Qwen3-0.6B \
    --samples 32 --steps 5 --batch-size 2 --block-size 4

# Smoke evaluation
python scripts/smoke_eval.py \
    --target Qwen/Qwen3-0.6B \
    --checkpoint /path/to/checkpoint.pt --gamma 7

# Benchmark comparison
python scripts/run_benchmark.py \
    --target Qwen/Qwen3-0.6B \
    --prompts examples/prompts.jsonl --gamma 7
\end{verbatim}

For GPU-based training at scale:
\begin{verbatim}
colab run --gpu A100 --keep --timeout 28800 \
    python scripts/smoke_train.py \
        --target Qwen/Qwen3-4B \
        --samples 10000 --steps 1000 --batch-size 8 --block-size 7
\end{verbatim}

% ======================================================================
\section{Conclusion}
\label{sec:conclusion}

We have presented \UraionSpec, a faithful, modular, and verified
implementation of the \DSpark{} speculative decoding algorithm. Our
implementation covers all core algorithmic components---semi-autoregressive
generation via Markov, gated Markov, and RNN sequential heads, confidence
prediction with analytical supervision, Sequential Temperature Scaling
calibration, and the hardware-aware prefix scheduler (Algorithm~1).
We also provide a standalone \DFlash-style parallel backbone module with
target model KV injection, available for integration into the draft model pipeline.

With 80 passing unit tests, comprehensive documentation, and reproducible
smoke training and evaluation pipelines, \UraionSpec provides a solid
foundation for academic research, ablation studies, and educational
exploration of speculative decoding. We hope this implementation serves as a
useful reference for researchers and practitioners working on efficient LLM
inference.

\paragraph{Future work.}
Key directions for extending \UraionSpec include full-scale GPU training on
the Open-PerfectBlend dataset, real engine throughput profiling for the
scheduler, multi-GPU distributed training via DeepSpeed or FSDP~(Fully Sharded Data Parallelism), integration
with production serving frameworks (vLLM, llama.cpp), and tree-based
verification for autoregressive drafters~\cite{spectrindecoding}.

% ======================================================================
\bibliographystyle{unsrt}
\begin{thebibliography}{20}

\bibitem{cheng2026dspark}
X.~Cheng, X.~Yu, C.~Shao, J.~Li, Y.~Xiong, et al.
\newblock DSpark: Confidence-scheduled speculative decoding with
semi-autoregressive generation.
\newblock \textit{ICML}, 2026.

\bibitem{leviathan2023fast}
Y.~Leviathan, M.~Kalen, and Y.~Matias.
\newblock Fast inference from transformers via speculative decoding.
\newblock \textit{ICML}, 2023.

\bibitem{chen2023accelerating}
C.~Chen, S.~Borgeaud, G.~Irving, J.-B.~Lespiau, L.~Sifre, and
J.~Jumper.
\newblock Accelerating large language model decoding with speculative
sampling.
\newblock \textit{arXiv:2302.01318}, 2023.

\bibitem{pope2023efficiently}
R.~Pope, S.~Douglas, A.~Chowdhery, J.~Devlin, J.~Bradbury, et al.
\newblock Efficiently scaling transformer inference.
\newblock \textit{MLSys}, 2023.

\bibitem{stern2018blockwise}
M.~Stern, N.~Shazeer, and J.~Uszkoreit.
\newblock Blockwise parallel decoding for deep neural machine translation.
\newblock \textit{arXiv:1811.03115}, 2018.

\bibitem{xia2024medusa}
H.~Xia, T.~Ge, S.-M.~Wang, S.-Q.~Chen, F.~Wei, and Z.-Y.~Shao.
\newblock Medusa: Simple LLM inference acceleration framework with multiple
decoding heads.
\newblock \textit{arXiv:2401.10774}, 2024.

\bibitem{miao2025eagle3}
Y.~Miao, Z.~Bai, Z.~Wang, J.~Zhou, and J.~Jia.
\newblock Eagle3: Efficient inference for LLMs.
\newblock 2025.

\bibitem{deepseekspec2026}
DeepSeek-AI.
\newblock DeepSpec: Production-grade speculative decoding for LLMs.
\newblock \url{https://github.com/deepseek-ai/DeepSpec}, 2026.

\bibitem{spectrindecoding}
Y.~Zhou, N.~Du, Z.~Zhong, T.~Ji, and Y.~Yang.
\newblock SpecTr: Tree-based speculative decoding.
\newblock \textit{EMNLP}, 2024.

\bibitem{dsparkalphaxiv}
DeepSeek-AI.
\newblock DSpark paper at alphaXiv.
\newblock \url{https://www.alphaxiv.org/abs/2026.dspark}, 2026.

\bibitem{uraionspechf}
Uraion Labs.
\newblock UraionSpec on HuggingFace.
\newblock \url{https://huggingface.co/UraionLabs/UraionSpec}, 2026.

\bibitem{uraionspecgh}
Arnav Prabhu.
\newblock UraionSpec on GitHub.
\newblock \url{https://github.com/arnavprabhu/UraionSpec}, 2026.

\end{thebibliography}

\end{document}
